\documentclass[11pt]{article}
\usepackage[margin=1in]{geometry}
\usepackage{mathptmx}
\usepackage{amsmath}
\usepackage{graphicx}
\usepackage{booktabs}
\usepackage{longtable}
\usepackage{setspace}
\usepackage[hidelinks]{hyperref}
\usepackage{caption}
\captionsetup{font=small,labelfont=bf}
\setstretch{1.15}
\usepackage{titlesec}
\titleformat{\section}{\normalsize\bfseries}{\thesection.}{0.5em}{\MakeUppercase}
\titleformat{\subsection}{\normalsize\bfseries}{\thesubsection}{0.6em}{}

\title{Bitcoin's Required Return}
\author{Agisilaos Papadogiannis\thanks{Independent Researcher (PhD, MBA). \emph{Correspondence:} apapadogiannis@gmail.com. \emph{ORCID:} 0009-0002-7753-7052. Data, code and replication files: \url{https://github.com/apapadogiannis/bitcoin-required-return}.}}
\date{\normalsize First version: November 6, 2026. This version: November 6, 2026.}

\begin{document}
\maketitle

\begin{abstract}
\noindent Bitcoin carries three incompatible ``expected returns'': a backtested drift of about 32\% a year (the age rule), a CAPM required return of 6--18\%, and the 14--22\% its holders say they expect. We reconcile them through Merton's (1987) limited-participation pricing, with every input measured rather than estimated. From a panel of 1,815 cryptoassets with 520 curated deaths, we estimate the first pricing-oriented mortality curves for the asset class: the hazard of death rises with age, rejecting the Lindy effect, and the catastrophic-death hazard of a top-20 asset is 0.3\% a year (LUNA, UST, FTT), capping the mortality premium at a few points. From household surveys and option prices we bound the required return of Bitcoin's marginal holder; realized returns then decompose into that required return plus an adoption wedge that was enormous early, zero by 2019, and negative after 2021. The age rule tracked the decaying sum, and its engine is spent. Expected return for 2026--2031: 9--22\% a year.
\end{abstract}

\medskip\noindent\textbf{\large Highlights}\par\nopagebreak
\begin{itemize}
\item From 1,815 cryptoassets and 520 curated deaths, the first mortality curves for the asset class built for pricing: the hazard of death rises with age, and the Lindy hypothesis behind maturity-premium stories is rejected.
\item Catastrophic death at the top is rare: among the 219 assets that ever reached the top twenty, the jump-collapse hazard is 0.30\% a year, so survival risk prices at a few points of expected return, not twenty-seven.
\item Bitcoin's required return from Merton's limited-participation condition with measured inputs (holders' surveyed 19\% portfolio weight, option-implied bounds) is 9--22\% a year; the age rule's 32\% tracked an adoption wedge that turned negative after 2021, and the paper commits to its downward break for September 2026--2031.
\end{itemize}

\medskip
\noindent\textbf{JEL classification:} G12, G11, G17, G41.\\
\textbf{Keywords:} Bitcoin, cryptoassets, required return, expected return, survival analysis, mortality, limited participation, marginal holder, Lindy effect, household finance.

\section{Introduction}

Suppose you must write down one number: the return Bitcoin should be expected to earn, per year, over the next five years. Three candidate answers are on the table, and they disagree by a factor of three.

The first comes from Bitcoin's own history. A simple rule (take the exponent of Bitcoin's power-law price trend, about 5.6 (Burger 2019; Papadogiannis 2026c), and divide it by Bitcoin's age in years) turned out, when tested as a forecast against the returns that actually followed, to be the only candidate that was right on average over 2012--2026 (Papadogiannis 2026c). Applied today it says about 32\% a year.\footnote{Units, once and for all: the age rule's instantaneous (arithmetic mean) drift at age $A$ is $\beta/A \approx 32\%$ at $A = 17.4$; its five-year-average form, used in the tables, is $(1+5/A)^{\beta/5}-1 \approx 33\%$; and its geometric (volatility-drag-adjusted) counterpart is 22--28\%. We quote 32\% as the headline and say explicitly, in Section 7, which form each test uses.} The second comes from standard asset pricing. The CAPM asks how much Bitcoin's risk would cost a diversified investor, and because Bitcoin barely moves with the stock market, the answer is modest: 6--18\% a year across the beta windows and premium assumptions surveyed in the companion paper (Papadogiannis 2026c); the 2024--26 consensus inputs used in Section 5 put it at 6--9\%. The third comes from asking people. In large household surveys, those who actually hold crypto said they expected 22\% over the next year in 2021 and 13.8\% in 2025 (Candia, Coibion, Gorodnichenko, and Weber 2026).

These numbers cannot all be the expected return. Yet each is produced by a method that is respectable on its own terms. The companion paper that documented the first number, the age-rule backtest of Papadogiannis (2026c), closed by posing exactly this puzzle: is the gap between what Bitcoin delivered and what any model says it should deliver \emph{compensation for risks that never materialized}, or \emph{the mark of an asset the market underpriced}? Realized returns alone cannot say, because Bitcoin is one asset with one history, studied precisely because it survived.

This paper answers the question by measuring the two things the time series cannot identify.

\textbf{First, we measure the risks that never materialized, by looking at the assets for which they did.} We build a panel of 1,815 cryptoassets from 2010 to 2026, including 520 that died, and estimate the asset class's mortality curves the way a demographer would. (Earlier descriptive studies of ``coin death'' exist, notably Feder, Gandal, Hamrick, and Moore (2018) and Grobys and Sapkota (2020); what is new here is a curated, pricing-oriented census: death types, dates, recovery fractions, size- and elite-conditional hazards, and the Lindy test.) Two results matter. The popular ``Lindy'' intuition, under which the longer an asset has survived the safer it is, with hazard falling like one over age, is rejected: the death rate of cryptoassets \emph{rises} with age. And the hazard that matters for Bitcoin, the death rate of assets that reached the top twenty, is small: 3.5\% a year for death by any cause, and only 0.30\% a year for the sudden, value-destroying collapses (Terra's LUNA and UST, and FTX's FTT, all in 2022). Priced at near-total loss, survival risk is worth a few points of expected return. It cannot explain a gap of roughly 27 points (the age rule's 32\% less the $\sim$5\% risk-free rate).

\textbf{Second, we bound the return Bitcoin's actual owners require, from their own portfolios.} The investor whose buying and selling sets Bitcoin's price is not the diversified investor of the CAPM. Surveys show the average crypto-owning household held 19\% of its financial assets in crypto in 2021; option prices from the same era imply the marginal trader was closer to all-in. For such an investor, elementary portfolio mathematics ties the return that holding Bitcoin must promise to three observable quantities: volatility, portfolio weight, and risk aversion. We call the resulting formula the marginal-holder condition (Section 2, equation~\ref{eq:foc}): risk aversion, times variance, times portfolio weight, plus small market and mortality terms. The route, in other words, is Merton's (1987) limited-participation pricing with every input measured, not a factor regression, not a consumption model, and not a cash-flow valuation, none of which can see the investor who actually clears this market. It is consistent with what holders said they expected in 2021 \emph{and} in 2025 at ordinary levels of risk aversion (1.7 and 2.8). The required return was genuinely enormous when Bitcoin's volatility was 120\% and its holders were concentrated, and it has collapsed, term by measurable term, to 9--22\% today.

With both measurements in hand, Bitcoin's history reads differently. Realized returns decompose into the required return of the marginal holder plus a \emph{wedge}: the part the holder was not promised and could not have demanded. The wedge was on the order of 100 points a year in 2012--2015, near zero by 2019, and negative after late 2021, when holders required 22--58\% and received about 14\%. The age rule, we show, is the smooth one-over-age curve drawn through required-plus-wedge; it backtested as unbiased because the two components happened to decay together over the sample, and it has drifted high since 2020 (a drift its own out-of-sample record already shows) because the wedge has run out and changed sign.

The paper's contributions, in order: (i) mortality curves for cryptoassets built for pricing, with the Lindy hypothesis tested and rejected and the top-asset hazard measured; (ii) a required-return estimate for Bitcoin built on Merton's (1987) limited-participation condition from its holders' measured portfolios, their stated expectations, and option-implied bounds: three sources, one equation; (iii) a decomposition of Bitcoin's realized drift that explains both why the age rule worked and why it has begun to fail; and (iv) a dated, falsifiable forecast: 9--22\% a year for September 2026--2031, against the age rule's 32\%.

Section 2 sets out the marginal-holder condition, in words first. Section 3 describes the data. Section 4 presents the graveyard. Section 5 measures what holders require. Section 6 rereads Bitcoin's history as required return plus wedge. Section 7 states the forecast and what would falsify it; Section 8 translates the results for an allocator; Section 9 lists limitations; Section 10 concludes.

\section{Two Returns, and the Person Who Sets the Price}

Everything in this paper runs through one equation, so we build it slowly and then display it.

\textbf{Expected versus required.} An asset's \emph{expected return} is a forecast: what it will earn on average. Its \emph{required return} is a hurdle: what an investor needs it to earn to be willing to hold it, given its risk. For a mature asset in a calm equilibrium the two coincide; the price adjusts until the forecast equals the hurdle. For an asset still being adopted, they can differ for years, and the difference is where portfolio decisions live (Papadogiannis 2026a, 2026b). The age rule is an estimate of the first object; the CAPM figures are estimates of the second; this paper's job is to measure the second properly and then see what remains of the first.

\textbf{Whose hurdle?} Here is the step most discussions skip. An asset has as many required returns as it has potential investors, but only one of them is reflected in the price: the hurdle of the \emph{marginal holder}, the investor who is just indifferent between holding a little more and a little less. If a pension fund would be happy to hold Bitcoin at an expected 8\%, but pension funds are not allowed to hold it, their 8\% is irrelevant to the price. The price is set by whoever actually clears the market.

\textbf{The marginal holder's arithmetic.} Take an investor with ordinary constant relative risk aversion, denoted $\gamma$ (empirically, $\gamma$ between 1 and 4 describes most people; 2--3 is the workhorse range). Suppose she holds a fraction $w$ of her wealth in an asset whose annual return volatility is $\sigma$, and whose return is essentially uncorrelated with the rest of her portfolio, a good description of Bitcoin, whose correlation with equities has lived below 0.2 for most of its life (Liu and Tsyvinski 2021). Her portfolio is optimal only if the asset's expected excess return equals the marginal risk it adds. That requirement is the central equation of the paper, which we call the \textbf{marginal-holder condition}:
\begin{equation}
\boxed{\;\mu - r_f \;=\; \underbrace{\beta\,\mathrm{ERP}}_{\text{market risk}} \;+\; \underbrace{\gamma\,\sigma^2\,w}_{\text{concentration}} \;+\; \underbrace{\lambda\,(1-\delta)}_{\text{mortality}}\;}
\label{eq:foc}
\end{equation}
Read the terms left to right. The left side is the asset's expected return over the risk-free rate. On the right: a small CAPM-style term for whatever market risk the asset does carry (beta times the equity risk premium); the \textbf{concentration term}, $\gamma$ times variance times the holder's own portfolio weight, which prices the risk she cannot diversify because she chooses, or is forced, to hold a lot of one thing; and a \textbf{mortality term}: if the asset can die suddenly with probability $\lambda$ per year, losing fraction $(1-\delta)$ of its value, the price must drift upward by $\lambda(1-\delta)$ in the surviving states just to break even. This last term is why a backtest run on a surviving asset measures something \emph{higher} than the unconditional expected return: the crashes that would have pulled the average down are missing from the sample by construction (Brown, Goetzmann, and Ross 1995).

\textbf{The mortality term, carefully.} The two symbols answer two different questions. $\lambda$ is the probability per year that the death event occurs. $\delta$ is the \emph{recovery fraction}: the share of her money a holder keeps when it does, so $(1-\delta)$ is the share she loses. They must be multiplied because dying and losing everything are not the same thing. When TerraUSD collapsed, holders kept about two cents on the dollar ($\delta \approx 0.02$). When BUSD was wound down under regulatory order, a death by any institutional definition, holders redeemed at par: $\delta \approx 1$, and nobody lost a cent. When a token merely migrates tickers, $\delta = 1$ exactly, which is why Appendix A refuses to count migrations as deaths at all. The product $\lambda(1-\delta)$ is therefore the \emph{expected loss rate from mortality}: probability per year, times fraction lost when it happens. Why it belongs in the drift: over a short interval $dt$, the holder earns the surviving drift $\mu\,dt$ with probability $1-\lambda\,dt$, and loses $(1-\delta)$ with probability $\lambda\,dt$, so her expected return counting both outcomes is $\mu\,dt - \lambda(1-\delta)\,dt$. Setting that equal to the required $r\,dt$ and rearranging gives $\mu = r + \lambda(1-\delta)$: the price must climb, while the asset lives, by exactly the expected death loss, an actuarially fair insurance premium paid out in drift, so that averaged over survival and death the holder earns what she requires. Numerically, a 3\% annual death probability with 95\% lost on death costs $0.03 \times 0.95 \approx 2.9$ points a year of drift. This is also the precise sense in which a survivor's backtest overstates the expected return: the backtest measures $\mu$, the drift on the paths where death did not happen; the unconditional expected return is $\mu - \lambda(1-\delta)$.

Equation~(\ref{eq:foc}) is Merton's (1987) economy of incomplete participation in its simplest clothing: when only a fraction of investors hold an asset, those who do must be paid for the idiosyncratic variance they absorb, and the payment scales with how much of it each must absorb. Every empirical section of the paper measures one of its terms: $\lambda(1-\delta)$ from the graveyard (Section 4), $\sigma$, $w$, and $\gamma$ from prices, surveys, and options (Section 5); Section 6 then uses the condition date by date, Section 7 uses it as a forecast, and Section 8 inverts it into a sizing rule.

\textbf{What equation~(\ref{eq:foc}) is, and is not.} Nothing in equation~(\ref{eq:foc}) is new, and that is deliberate. The concentration term is the first-order condition of textbook portfolio choice (Merton 1969), and its use as a pricing statement under limited participation is Merton (1987); the mortality term is the standard jump compensation (Merton 1976) read through the survival conditioning of Brown, Goetzmann, and Ross (1995); the beta term is the CAPM. The paper's contribution is not the equation but its inputs. Three clarifications follow. First, equation~(\ref{eq:foc}) is an optimality condition that any equilibrium must satisfy, not an equilibrium model: we do not model supply, endogenize $w$, $\sigma$, or $\lambda$, or derive the price level. It therefore complements, rather than competes with, the equilibrium models of Bitcoin's price (Biais, Bisi\`ere, Bouvard, Casamatta, and Menkveld 2023; Pagnotta 2022; Cong, Li, and Wang 2021), supplying the discount-rate input those models leave to be estimated. Second, the condition is generic by design: the same line prices founder stock, private equity, and emerging-market equities before liberalization. Bitcoin enters through the values of the inputs, not the functional form, and through one data advantage no other asset class offers: a thousand-asset graveyard from which $\lambda$ can be measured rather than assumed. The emerging-market integration literature (Bekaert and Harvey 2000; Henry 2000) documents exactly the mechanism our forecast rests on: required returns are high while a market is segmented and compress when it opens to outside capital; the ETF era is Bitcoin's liberalization. Third, using the survey-average holder weight for $w$ is not casual: in an equilibrium of holders with similar risk aversion, every unconstrained holder's condition binds at the same $\gamma\sigma^2 w$, so weights equalize across participants and the average participant's weight identifies the marginal one's.

\textbf{A worked example.} In 2021, Bitcoin's volatility was about 80\% a year, so $\sigma^2 = 0.64$. Survey data (Section 5) show the average crypto-owning household held $w = 19\%$ of its financial assets in crypto. With $\gamma = 2$ the concentration term is $2 \times 0.64 \times 0.19 \approx 24\%$ a year. Add a risk-free rate near 1.5\%, a couple of points of market beta, and a fraction of a point of mortality, and the marginal holder of 2021 required roughly 25--30\% a year. This is why, as we will see, holders \emph{saying} they expected 22\% were not being exuberant. They were reading their own first-order condition.

\textbf{What makes each term fall.} Volatility falls as markets deepen: Bitcoin's $\sigma$ has gone from about 150\% (2013) to about 40--50\% (2026). The weight $w$ falls as participation broadens: when the only people willing to hold an asset are enthusiasts, each must hold a lot of it; when ETFs let everyone hold a sliver, nobody has to. And $\lambda$ falls, or at least the market's estimate of it falls, as the asset survives attacks it was feared not to. Each term is measurable, each has fallen, and the sum today is a single-to-low-double-digit number.

\section{Data}

The asset panel comes from the Coin Metrics community dataset (Coin Metrics 2026), the same source the companion backtest paper uses for Bitcoin and Ether prices, observed from July 2010 to May 24, 2026. After removing fiat exchange rates, tokenized equities, wrapped and staked duplicates, stablecoins, and trading-pair artifacts, the core universe is \textbf{1,815 cryptoassets}, contributing about 67,700 asset-months of exposure. The dataset carries full daily \emph{trading volume} for essentially all of them, and full daily \emph{prices} for the 139 historically significant assets; volume therefore carries the life-and-death timeline, and prices carry the return calculations on the significant subset.

\textbf{Birth.} An asset is born, for our purposes, on the first day it shows a price or trading volume: the day it enters the investable, tracked universe. This is the clock an investor cares about.

\textbf{Death.} Death is harder, because cryptoassets rarely die cleanly. We use three mechanical definitions and then a curated layer:
\begin{enumerate}
\item \textbf{Feed death}: the data feed ends more than 180 days before the dataset does. This proxies delisting from tracked markets. 234 cases.
\item \textbf{Price death}: the price falls below 1\% of its all-time high and never recovers (measurable only on the price-complete subset). 33 cases.
\item \textbf{Volume death}: trailing three-month volume falls below 1\% of its historical peak \emph{and} below an absolute floor of \$250{,}000 a day, and never recovers. The absolute floor matters: without it, healthy mid-caps whose volume merely normalized after the 2021 mania would be falsely counted as dead. 438 cases, overlapping the other two definitions.
\item \textbf{Curation.} Two corrections are applied by hand, with the full audit trail in Appendix A. \emph{Ticker migrations are not deaths}: when MATIC became POL, EOS became Vaulta, or Golem's GNT became GLM, holders received successor tokens at full value, so these events (seventeen in all) are removed. And \emph{known catastrophes are dated correctly}: the free data tier hides the price history of several major assets, so the collapses of LUNA and UST are entered at May 12, 2022 and FTT at November 10, 2022 from the public record, rather than at whatever date a volume rule would eventually notice them.
\end{enumerate}

The result: \textbf{520 deaths among 1,815 assets (29\%)}, of which 8 are classified as sudden, value-destroying \emph{jump-collapses} and the rest as gradual abandonment: projects that faded, lost their listings, and stopped trading.

The other inputs are: Bitcoin's monthly price, market capitalization, and count of addresses with a balance (Coin Metrics); US household crypto-ownership rates from the Federal Reserve's payment and household surveys and Pew (0.4\% of adults in 2014, 1.2\% in 2018, 9--12\% in 2021, roughly 11--14\% since); holders' portfolio shares and expectations from Candia, Coibion, Gorodnichenko, and Weber (2026), with Austrian household-finance data (OeNB 2025) as a European cross-check; ten-year Treasury yields; global household wealth from the UBS Global Wealth Report; and option-implied figures from Deribit as reported in the literature and the DVOL index.

One bias deserves a sentence up front. Coin Metrics adds an asset to its universe only once the asset has become significant, so thousands of micro-cap tokens that died young and obscure never enter the panel. This \emph{understates} the death rate at young ages, which means it biases the test \emph{against} our headline finding that hazard rises with age. The elite estimates, which are what Bitcoin's pricing needs, are unaffected: once an asset is in the top twenty it is fully observed from then on.

\section{The Graveyard: How Cryptoassets Die}

A \emph{hazard rate} answers a simple question: of the assets alive at a given age, what fraction die within the next year? If 100 three-year-old coins are alive in January and 12 are dead by December, the hazard at age three is 12\% a year. Estimating hazards across ages, with assets that are still alive treated correctly as ``censored'' (we know they lived \emph{at least} this long), is standard survival analysis; the machinery is the same used for patients, firms, and light bulbs.

\textbf{The Lindy question.} The age rule's most natural economic defense was mortality: perhaps Bitcoin's enormous drift was fair payment for a risk of sudden death, and perhaps that risk declines as the asset ages. This is the ``Lindy effect,'' under which a technology's remaining life expectancy grows with its age, so the hazard falls like $k$ divided by age. If that were true with $k$ near 5.6, the age rule would simply be survival compensation, rationally priced, with equation~(\ref{eq:foc})'s mortality term doing all the work. The hypothesis is attractive, parameter-free in shape, and testable; testing it is the point of building the graveyard.

\textbf{It fails.} Table~\ref{tab:hazard} shows the death rate by age bucket under the curated definition.

\begin{table}[ht]\centering\small
\caption{Mortality hazard by age, curated death definition (1,815 assets, 520 deaths, 2010--2026).}
\label{tab:hazard}
\begin{tabular}{lccc}
\toprule
Age (years) & Deaths & Exposure (asset-yrs) & Hazard (per yr) \\
\midrule
0.1--0.5 & 6 & 693 & 0.9\% \\
0.5--1 & 37 & 770 & 4.8\% \\
1--2 & 111 & 1{,}213 & 9.2\% \\
2--3 & 106 & 905 & 11.7\% \\
3--4 & 82 & 661 & 12.4\% \\
4--5 & 78 & 488 & 16.0\% \\
5--7 & 63 & 502 & 12.6\% \\
7--9 & 33 & 198 & 16.7\% \\
9--12 & 4 & 53 & 7.5\% \\
\bottomrule
\end{tabular}
\par\smallskip\parbox{0.8\textwidth}{\footnotesize Buckets begin at age 0.1 (the first weeks of a listing are definitionally alive) and the sparse 12--17 bucket (7 asset-years, no deaths, Bitcoin the main occupant) is omitted; total exposure across all ages is about 5,640 asset-years.}
\end{table}

The hazard \emph{rises} with age through at least age nine. Formally, a Poisson regression of monthly deaths on log age, with calendar-year fixed effects to absorb crypto winters and standard errors clustered by asset, gives an age elasticity of \textbf{+0.42 (s.e.\ 0.04)}: each doubling of age \emph{raises} the death rate by roughly a third, holding the calendar year fixed (the fixed effects absorb market-wide waves such as crypto winters; clustering lets each asset's months be correlated). Lindy requires $-1$. The estimate has the wrong sign by eleven standard errors, and it is robust: $+0.44$ to $+0.56$ across the three mechanical death definitions, $+0.65$ when the suspicious cluster of mid-2025 feed terminations is excluded, and positive within the elite subsample alone. Equivalently, the slope of log survival against log age, which equals $k$ if the hazard really were $k/\text{age}$, gives \textbf{between about 0.2 and 0.5} depending on the death definition, \textbf{against the 5.6} the age rule would need. A 5.6 would have implied that a three-year-old cryptoasset had roughly a 1-in-20,000 chance of reaching Bitcoin's current age; the measured elite hazard, held constant, implies a 14-year survival probability around 0.6. Cryptoassets age like small firms, with a liability of adolescence and obsolescence, not like Lindy classics (Figure~\ref{fig:hazard}).

\textbf{Where the data end, honestly.} Three support caveats belong next to that result. First, in the oldest populated bucket (ages 9--12) the raw hazard falls back to 7.5\%, on 4 deaths and 53 asset-years; the rise in hazard is established through roughly age nine, the data's support ends near age twelve, and Bitcoin is 17.4. The Lindy rejection is therefore a statement about the asset class over its observed life span, not a measurement of Bitcoin's own current hazard; that is exactly why the pricing below uses the elite hazard \emph{level} rather than any age-extrapolation. Second, the three jump-collapses sit within six months of one another and were causally linked (the Terra collapse propagated through Celsius to FTX), so the Poisson confidence interval, which assumes independent events, understates the true uncertainty; we treat 0.30\% as the point estimate of a partly systematic process and charge for that in Section 8. Third, the 2024--25 delisting purge supplies 12 of the 36 elite deaths; treating every 2025 death as censoring instead of death lowers the all-cause elite hazard from 3.5\% to 2.4\% a year (CI 1.5--3.5\%), so the mortality bound of $\leq$3 points used below is conservative under either reading.

\begin{figure}[ht]\centering
\includegraphics[width=0.78\textwidth]{fig1_hazard_age.pdf}
\caption{Mortality hazard versus age (log--log). The fitted line uses the period-fixed-effects elasticity of $+0.42$; the dashed line is the Lindy shape $\lambda=k/A$.}
\label{fig:hazard}
\end{figure}

\textbf{Size is what protects, not age.} Sorting assets by peak trading volume (Table~\ref{tab:size}), the hazard falls monotonically and steeply with size.

\begin{table}[ht]\centering\small
\caption{Mortality by size and for elite cohorts (curated deaths).}
\label{tab:size}
\begin{tabular}{lcc}
\toprule
Size group & $n$ & Hazard (per yr) \\
\midrule
Smallest quartile & 452 & 17.4\% \\
Q2 & 451 & 13.0\% \\
Q3 & 451 & 9.8\% \\
Largest quartile & 452 & 3.2\% \\
Ever top-50 by volume & 499 & 5.5\% (CI 4.5--6.5\%) \\
\textbf{Ever top-20 by volume} & \textbf{219} & \textbf{3.5\% (CI 2.5--4.9\%)} \\
Ever top-20, jump-collapse only & 219 & \textbf{0.30\% (CI 0.06--0.86\%)} \\
\bottomrule
\end{tabular}
\end{table}

The last two lines are the ones Bitcoin's pricing needs. Among the 219 assets that ever reached the top twenty by monthly volume, 36 died over 1,017 asset-years of exposure, an all-cause rate of 3.5\% a year, and only three of those deaths were the sudden, catastrophic kind that a jump-risk premium prices: LUNA and UST in May 2022 and FTT in November 2022, each a roughly total loss within days. None has occurred since. Excluding four contested cases, assets flagged dead by the drawdown rule but whose chains still operate (ICP, SUSHI, FLOW, NANO; see Appendix A), the all-cause rate falls to 3.15\% (CI 2.2--4.4\%). The remaining elite deaths are slow fades: OMG, NEM, Bitcoin Gold and their like lost their listings over months or years, during which holders could exit at declining prices. Gradual losses of that kind live in an asset's volatility and drift, not in its jump hazard. The 520 dead assets are, in the taxonomy of the companion scarcity framework (Papadogiannis 2025), the ``physically scarce but no longer desired'' quadrant realized at scale: hard-coded supply caps outlived demand, and scarcity without desirability monetizes nothing (Figure~\ref{fig:size}; the survival curves by size are Figure~\ref{fig:km}; the death timeline, with the 2022 cascade and the 2024--25 delisting purge visible, is Figure~\ref{fig:timeline}).

\begin{figure}[ht]\centering
\includegraphics[width=0.78\textwidth]{fig2_size.pdf}
\caption{Hazard by peak-volume quartile and for the elite cohorts; whiskers are exact Poisson 95\% intervals.}
\label{fig:size}
\end{figure}

\begin{figure}[ht]\centering
\includegraphics[width=0.80\textwidth]{fig5_km_size.pdf}
\caption{Kaplan--Meier survival curves by peak-volume quartile, with the ever-top-20 cohort (dashed). Half of the smallest quartile is dead by age 3.3; the largest quartile and the elite track each other and remain above 60\% survival past age 10.}
\label{fig:km}
\end{figure}

\begin{figure}[ht]\centering
\includegraphics[width=0.78\textwidth]{fig4_timeline.pdf}
\caption{When the elite died: 36 deaths among 219 ever-top-20 assets, by year and type.}
\label{fig:timeline}
\end{figure}

\textbf{What this buys the pricing equation.} Price the jump channel at near-total loss: $\lambda(1-\delta) \approx 0.3\% \times 0.95 \approx$ \textbf{0.3 points a year}. Treat every elite death, even the slow fades, as if it were a jump at total loss and you still cannot get above $\approx$ \textbf{3 points}. The mortality term in equation~(\ref{eq:foc}) is real (2022 proved it can strike the top ten), but it is a few points, not twenty-seven. One caveat is honest and important: no cross-section can measure Bitcoin's \emph{own} hazard in 2010--2014, when it was the only asset of its kind and plausibly faced genuinely existential risk; our bounds apply from the point at which an elite class exists to measure. Section 6 keeps that early window wide.

A by-product worth recording: the deaths cluster in \emph{calendar time} (2022, 2024--25), not smoothly in age. Mortality at the top arrives with contagion episodes and delisting purges. To the extent the hazard is systematic rather than idiosyncratic, a diversified allocator may rationally charge for part of it; we return to this in Section 8.

\textbf{Forward-looking hazards the graveyard cannot see: the quantum threat.} The hazard measured here is a backward-looking frequency, and some risks have never yet fired for anyone. The clearest is quantum computing: a cryptographically relevant quantum computer running Shor's algorithm would break the elliptic-curve signatures that secure Bitcoin, and the event would be exactly what the marginal-holder condition calls a jump, a $\lambda$-event with its own loss fraction. The model therefore \emph{covers} the quantum threat by construction; what the graveyard cannot do is price it, because the realized death count contains no quantum events. Two observations discipline it anyway. First, the loss given the event is scenario-dependent, not automatically total: roughly a quarter of Bitcoin's supply sits in addresses with exposed public keys, including the presumed-Satoshi coins (Deloitte's standing estimate is about 4 million BTC) and is the immediately vulnerable part, while a telegraphed advance in quantum hardware would allow a migration to post-quantum signatures that preserves most other value. So $\delta$ ranges from near zero (a surprise break) to substantial (an orderly transition). Second, if the market prices a material quantum hazard, it must appear where all jump risk appears: in deep out-of-the-money put prices. The option-implied crash intensity is therefore the running referee between the graveyard's 0.3\% and any larger forward-looking hazard, quantum included; at this writing it shows no such wedge. An investor who privately assigns, say, a 1\% annual probability to a disorderly quantum break with half of value lost should simply add 0.5 points to the mortality term; the equation accepts the view without modification.

\section{What Bitcoin's Holders Require}

The marginal-holder condition, equation~(\ref{eq:foc}), has four inputs on the right-hand side. The graveyard supplied $\lambda$. This section supplies the rest (volatility, the holder's portfolio weight, and risk aversion) from three independent sources, and shows they tell one story.

\textbf{Source 1: holders' measured portfolios.} The Cleveland Fed study of cryptocurrencies in household finance (Candia, Coibion, Gorodnichenko, and Weber 2026) reports the number the equation needs most: among US households that owned any crypto in 2021Q3, crypto made up \textbf{19\% of financial assets on average}, with 40\% of owners below 5\% and about a fifth of owner-households above \emph{half} their portfolio. A European cross-check points the same way: in the Austrian central bank's 2023 household survey, owners mostly held under 20\% of financial assets in crypto, rarely above 20--25\% (OeNB 2025). So the marginal holder of the early 2020s was not the diversified investor of the CAPM; she held a fifth of her wealth in the asset.

\textbf{Source 2: holders' stated expectations, a consistency check.} The same surveys asked holders what return they expected over the next year: 22\% in 2021, 13.8\% in 2025. Now run equation~(\ref{eq:foc}) forward at each date and ask what risk aversion $\gamma$ makes it hold (Table~\ref{tab:gamma}).

\begin{table}[ht]\centering\small
\caption{Two survey waves, one equation.}
\label{tab:gamma}
\begin{tabular}{lcc}
\toprule
 & 2021Q3 & 2025 \\
\midrule
Volatility $\sigma$ (trailing) & $\sim$80\% & $\sim$45\% \\
Holder portfolio weight $w$ & 0.19 (measured) & 0.08--0.19 (bounded) \\
Risk-free + beta term & $\sim$3--4\% & $\sim$7--8\% \\
Stated expected return & 22\% & 13.8\% \\
\textbf{Implied} $\boldsymbol{\gamma}$ & $\boldsymbol{\approx 1.7}$ & $\boldsymbol{\approx 2.8}$ \\
\bottomrule
\end{tabular}
\par\smallskip\parbox{0.86\textwidth}{\footnotesize Implied $\gamma$ solves equation~(\ref{eq:foc}) at $\sigma=0.77$, $w=0.19$, $r_f+\beta\mathrm{ERP}=3.5\%$ (2021) and $\sigma=0.45$, $w=0.11$, $r_f+\beta\mathrm{ERP}=7.5\%$ (2025), with the mortality term at 0.3 points. Across the stated input ranges the 2025 figure spans roughly 1.6--3.1.}
\end{table}

Two clarifications about what this table does and does not show. With one free $\gamma$ per wave, each line can be made to hold by construction; this is a range-consistency check, not a formal overidentification test. Its content is that the implied $\gamma$ lands both times inside the narrow interval (roughly 1 to 4) that describes ordinary investors, rather than at the implausible values (0.2, or 15) that would expose the condition as irrelevant to how holders actually think. And the move from 1.7 to 2.8 is itself economically sensible rather than an embarrassment: as participation broadens from early enthusiasts toward the general investing public, the marginal holder should become more risk-averse; with one $\gamma$ per wave this is a reading, not a test, and we flag it as such. Holders' expectations, so often dismissed as exuberance, are at least consistent with their own portfolio first-order condition, the same beliefs-to-portfolios coherence documented for equity investors by Giglio, Maggiori, Stroebel, and Utkus (2021). This is the paper's quiet central exhibit: the 14--22\% that ``holders say they expect'' was never a third, inconsistent number. It is the required return of the people who actually hold the asset, moving as its determinants move.

\textbf{Source 3: option prices.} Options give a holder-free reading. A classic result (Martin 2017) says the risk-neutral variance priced into options is a lower bound on the expected excess return demanded by an investor fully invested in the asset; in equation~(\ref{eq:foc}) the concentration term is $\gamma \sigma^2 w$, so ``$\gamma w = 1$'' describes a fully invested log-utility investor, whose required excess return is exactly the return variance; Martin's result is that option prices put a floor at that variance. (``Risk-neutral variance'' is the variance computed from option prices rather than from history.) Deribit-based estimates for 2017--2022 put that bound at \textbf{66--84\% a year}, with Bitcoin implied volatility then running at 80--90\% (``Risk Premia in the Bitcoin Market,'' 2024; sample July 2017--December 2022). This is the concentration mechanism priced by the market itself, during the era when crypto markets were documentably segmented from mainstream capital, with large and persistent cross-market price gaps (Makarov and Schoar 2020), and the marginal trader really was close to all-in. Today Deribit's DVOL index (the 30-day option-implied volatility, a VIX for Bitcoin) has run between about 40 and 55, so the same all-in bound is $0.40^2$--$0.55^2 \approx$ \textbf{16--30\%}, and the actual marginal holder, who in the ETF era is no longer all-in but holds $w$ of 8--19\%, requires $\gamma\sigma^2 w \approx$ \textbf{2--9 points} above the risk-free-plus-beta floor.

\textbf{Putting the 2026 number together.} Risk-free 5.0\%; beta term 1.3--4.2 points (beta 0.33--1.05 on a 4\% equity premium); concentration 2--9 points ($\gamma \in [1.7,2.8]$, $\sigma = 42\%$, $w \in [0.08,0.19]$); mortality 0.3--3 points. Summing: 8.6\% to 21.2\% a year, which we quote as \textbf{9--22\%} throughout, with a center near 13--15\%. For contrast, the age rule says 32\%. The two cannot both be the expected return, unless the difference is something other than a required premium. Section 6 shows that historically it was exactly that: a wedge.

\section{Rereading Bitcoin's History: Required Return Plus Wedge}

Write each period's realized return as the sum of two parts:
\begin{equation}
\text{realized} \;=\; \underbrace{\mu^{req}(t)}_{\text{what holders could demand}} \;+\; \underbrace{\varepsilon(t)}_{\text{the wedge: what nobody was promised}}.
\end{equation}
The required part, $\mu^{req}(t)$, is equation~(\ref{eq:foc}) evaluated date by date from the measured ingredients: trailing volatility from the price series; the holder weight $w$ bounded between a market-clearing back-cast and the all-in trader the option bounds describe; the measured mortality term. The wedge is what is left. The back-cast deserves a plain-words explanation, since it generates the whole pre-2021 lower band. Market clearing requires that Bitcoin's market capitalization equal the wealth its holders have placed in it: cap $= \bar{w} \times$ (participating wealth). We anchor participating wealth in 2021Q3, the one date where both sides are measured (capitalization divided by the surveyed 19\%), and extend it backward and forward in proportion to the number of participants (the survey-anchored adoption series of Section 3) times average wealth per participant (a ramp from a quarter of its 2021 level in 2012, when holders were students and hobbyists, to 1.4 times by 2026, when they include institutions). Dividing capitalization by that path gives the implied average holder weight at each date: about 2\% in 2013. Read literally, it says that \emph{if} early Bitcoin wealth had been spread across survey-style households, each held a small position and required little; the all-in alternative reads the same capitalization as concentrated in few fully committed hands who required a great deal. The truth sits between those readings, and the band is exactly that interval. Table~\ref{tab:wedge} does this at two-year intervals; Figure~\ref{fig:wedge} shows the full monthly series.

\begin{table}[ht]\centering\small
\caption{Required return versus realized drift versus the age rule.}
\label{tab:wedge}
\begin{tabular}{lccccc}
\toprule
Year & Age & Volatility & Required (band) & Realized next-5y CAGR & Age rule (5-yr form) \\
\midrule
2013 & 4.5 & 151\% & 10--231\% & 135\% & 132\% \\
2015 & 6.5 & 115\% & 7--135\% & 103\% & 90\% \\
2017 & 8.5 & 80\% & 7--67\% & 51\% & 68\% \\
2019 & 10.5 & 90\% & 14--82\% & 42\% & 55\% \\
2021 & 12.5 & 83\% & 22--58\% & $\sim$14\% (to 2026) & 46\% \\
2023 & 14.5 & 66\% & 12--23\% & -- & 40\% \\
2026 & 17.4 & 42\% & 9--22\% & -- & 33\% \\
\bottomrule
\end{tabular}
\end{table}

\begin{figure}[ht]\centering
\includegraphics[width=0.92\textwidth]{fig_wedge.pdf}
\caption{Required return versus realized forward five-year drift versus the age rule. The shaded gap between realized and required is the wedge: positive early, negative after late 2021. Green diamonds mark holders' surveyed expectations; the dotted line marks spot-ETF approval (January 2024), Bitcoin's market-integration event.}
\label{fig:wedge}
\end{figure}

Three readings of this table, in increasing order of importance.

\textbf{Why the bands are wide early.} The honest uncertainty in 2013--2017 is \emph{who} the marginal holder was. Under the survey-anchored back-cast, participants were few but each held a modest weight, and the required return was only 10--15\%: on that reading the wedge was enormous, 100-plus points a year; the market simply did not anticipate the adoption that came, and the realized returns were a surprise, not a payment. (In the language of the companion scarcity framework (Papadogiannis 2025), where network value is desirability times scarcity amplified by adoption, the wedge is the unanticipated growth of that value function; this paper prices the risk side, that framework describes the value side.) Under the all-in reading, which the 2017--22 option bounds support for the segmented era, required returns were themselves in the tens-to-hundreds, and the wedge was modest. The truth sits between, and for the paper's purpose it does not need to be settled: on \emph{either} reading the early engine (concentration, volatility, segmentation, surprise) has measurably shut down.

\textbf{Why the age rule backtested so well, and why it has started to miss.} Both the required return and the wedge decayed over 2012--2021, roughly together, and a one-parameter power law in age is an excellent smoother of a sum of decaying curves. That is the age rule: a reduced form of required-plus-wedge, not a structure. The reduced form has no engine of its own, so when the components stopped decaying \emph{together}, it started to miss, and in the one direction the decomposition predicts. From late 2021, holders required 22--58\% (their own first-order condition; the option bound was higher still) and received about 14\%: \textbf{the wedge turned negative.} The companion backtest paper's own out-of-sample record shows precisely this: in windows beginning after the power law's 2019 publication, Bitcoin fell short of the age rule by 5--8 points a year. What looked there like a small, possibly transient overshoot is, in this decomposition, the sign flip of the wedge.

\textbf{Why the wedge is unlikely to be a hidden risk premium.} Three facts narrow that escape route. The graveyard caps the mortality channel at $\sim$0.3--3 points (Section 4), and rejects the declining-hazard shape outright. The $\sigma^2 w$ concentration channel is measured, and today contributes 2--9 points. And, in a test we ran against ourselves, the concentration quantity has \emph{no forecasting power of the right sign}: regressing forward 3--5 year returns on real-time $\sigma^2 w$ alongside the age variable, $\sigma^2 w$ enters negative (because market cap sits in $w$'s numerator, the product spikes at mania tops) while the age term is undisturbed. The marginal-holder condition is an equilibrium restriction holding at each date; it is not, and the paper does not claim it to be, a timing signal. What these facts cannot exclude are premia for risks outside both measured channels, such as ambiguity aversion or disaster states never realized in 2010--2026; if present, they must show up where all priced tail risk shows up, in option prices, and Section 7 sends the reader there. Full regressions in Appendix B.

\textbf{The cross-section tells the same story from the side.} If early crypto returns had been mostly required premium, the typical altcoin (more volatile, more concentrated, higher $\lambda$ than Bitcoin) should have delivered even more. It delivered less than nothing: across the significant-asset cohorts, the \emph{median} coin born 2015--2020 shows a negative compound return at every horizon, while cohort \emph{means} stay large on the back of a few winners, exactly the lottery shape that a 10--17\% annual death rate plus a handful of extreme survivors produces. Within Bitcoin's own 2010--14 birth cohort at age ten, the mean coin compounded at about +167\% a year, the median at +50\%, and Bitcoin at +233\%. Conditioning on the winner, which every rule fitted to Bitcoin's history silently does, is worth on the order of 180 points a year over the in-cohort median at that age. The wedge, in other words, is visible twice: over time in Bitcoin's own decomposition, and across assets as the gap between the winner and the median.

\section{The Forecast, and What Would Prove It Wrong}

If the wedge is spent (no systematic adoption surprise left to harvest, the winner already identified), then the expected return going forward is just the required return of the marginal-holder condition, and every ingredient is now observable:

\textbf{Bitcoin's expected return for September 2026 to September 2031 is 9--22\% a year, centered near 13--15\%.} The age rule says 32\% (mean-type), 22--28\% compounded. The two forecasts differ enough to be separated by a single five-year window, and we commit to ours in print.

To keep the claim honest, here is what each outcome would mean:
\begin{itemize}
\item \textbf{BTC compounds (CAGR) at $\geq\sim$25\% a year to 2031.} The age rule is vindicated and this paper's decomposition is wrong in an instructive way: either a required-return component we measure as small is in fact large (which the options market would have to show; watch whether implied variance and holder concentration stop compressing), or systematic adoption surprises were not spent after all.
\item \textbf{BTC compounds in the single digits to low twenties.} The decomposition stands; the age rule's first structural break is confirmed, in the direction and for the reason stated here.
\item \textbf{BTC compounds near zero or below.} Within our band's logic this is a bad draw of the (large) diffusive risk, not a refutation. But if it comes with \emph{rising} option-implied crash intensity well above the graveyard's 0.3--3\%, the mortality term was underestimated.
\end{itemize}

Two honesty notes on the test itself. \emph{Units:} our 9--22\% is an arithmetic-drift band; at today's $\sigma \approx 42$--45\% the volatility drag is roughly $\sigma^2/2 \approx 9$ points, so the compounded (CAGR) counterpart is about 0--13\%, against the age rule's compounded 22--28\%. In matching units the two forecasts do not overlap, and a single five-year window grades them. \emph{Asymmetry:} the forecast is sharply falsifiable upward (a $\geq$25\% CAGR vindicates the age rule) but only conditionally falsifiable downward: within the band's logic a near-zero outcome is a bad draw of large diffusive risk unless it arrives with rising option-implied crash intensity. A one-and-a-half-sided test is the honest price of forecasting a 45\%-volatility asset, and we state it rather than hide it.

Two auxiliary predictions sharpen the mechanism. First, drift compression should be fastest where volatility and holder concentration fall fastest: Ether post-ETF should show the same compression; a still-segmented, still-concentrated large asset should not. Figure~\ref{fig:eth} shows the first check already passes: Ether's volatility and forward returns compress alongside Bitcoin's into the ETF era. Second, option-implied crash intensity should remain in the low single digits, consistent with the measured elite jump hazard; a reading persistently above $\sim$10\% a year would push our band up and signal that the market disagrees with the graveyard, whether about contagion, regulation, or forward-looking tail risks such as a quantum break of Bitcoin's signature scheme (Section 4), which realized mortality cannot yet contain.

\begin{figure}[ht]\centering
\includegraphics[width=0.95\textwidth]{fig6_eth.pdf}
\caption{The compression mechanism is not Bitcoin-specific. Left: EWMA annualized volatility for Bitcoin and Ether, with spot-ETF approval dates marked. Right: realized forward three-year returns for both assets.}
\label{fig:eth}
\end{figure}

\section{For the Allocator}

The decomposition translates directly into practice.

\textbf{The hurdle.} A diversified allocator should require of Bitcoin what its diversifiable-risk-adjusted contribution costs her: the risk-free rate plus beta times the equity premium, about \textbf{7--9\% today}, plus some or all of the mortality term, 0.3--3 points, to the degree she judges crypto mortality systematic rather than diversifiable. The 2022 clustering of elite deaths (Terra, Celsius, FTX within seven months) argues it is partly systematic; we would charge 1--2 points. Call the hurdle 8--11\%. This is not a forecast of what Bitcoin will do; it is the bar it must clear to belong in her portfolio at more than market weight.

\textbf{The sizing rule.} The marginal-holder condition, inverted, sizes the position. If an investor's own view of Bitcoin's expected return is $\hat{\mu}$, her risk aversion $\gamma$, and Bitcoin's volatility $\sigma$, her weight should be
\begin{equation}
w^{*} \;=\; \frac{\hat{\mu} - r_f - \beta\,\mathrm{ERP} - \lambda(1-\delta)}{\gamma\,\sigma^{2}}.
\end{equation}
Worked example: an investor with $\gamma = 3$ who believes our central forecast, $\hat{\mu} = 14\%$, at $\sigma = 45\%$: $w^{*} = (14 - 5 - 2.7 - 0.3)/(3 \times 0.2) \approx$ \textbf{10\% of financial wealth}. That is aggressive by institutional standards, and roughly half the 19\% the average 2021 holder actually ran at far higher volatility. An investor who believes the age rule's 32\% gets $w^{*} \approx 40\%$; the distance between those two positions is the practical cost of getting this paper's question wrong.

\textbf{The Black--Litterman translation.} In a Black--Litterman portfolio, the tilt away from market weight is proportional to the gap between your expected return and the market-implied one (Black and Litterman 1992). This paper reduces that gap to observable disagreements: your volatility path, your view of holder-base integration (the path of $w$), and your mortality estimate against the market's option-implied one. An allocator with no disagreement on those three has no basis for holding Bitcoin away from its market weight, currently well under 1\% of global financial wealth, however glorious its history.

\section{Limitations and the Measurement Agenda}

This is a measurement paper, so its limits are measurement limits; we list them in the order a critical reader should weigh them.
\begin{enumerate}
\item \textbf{The holder weight before 2021 is bounded, not measured.} The portfolio-share surveys begin in 2021; everything earlier rests on market-clearing back-casts and option bounds. The paper's historical bands are honest about this width, and its forward claim does not depend on the early years. The single most valuable upgrade would be a measured $w$ path: entity-adjusted on-chain concentration joined to holder-wealth estimates, the 2025 portfolio-share figures in the Candia et al.\ online appendix (not yet published in text at this writing), and the ETF share of total holdings.
\item \textbf{Coverage-onset selection in the graveyard.} Assets enter the panel only after becoming significant, so young-age hazards are understated. This biases the data \emph{toward} Lindy, which makes the rejection conservative, and it leaves the elite estimates untouched; but population hazard \emph{levels} at ages under two should be read as floors.
\item \textbf{Death timing and classification.} Feed and volume deaths lag economic death by months; four elite ``deaths'' are contested zombies (chains that operate far below their manic highs); one (ICP) is an artifact of a listing-day price spike. All classifications, dates, recovery fractions, and confidence flags are in Appendix A, and the headline elite hazard moves only from 3.5\% to 3.15\% when the contested cases are dropped.
\item \textbf{Bitcoin's own existential era is unmeasurable.} No cross-section existed in 2010--2014 to price the risk that Bitcoin itself was a dead end. Our early-period required-return bands are therefore wide by construction, and nothing in the paper's forward claim rests on resolving them.
\item \textbf{Survey-to-wealth pairing.} The participation series is US-based and the wealth aggregates global; re-running the clearing calculations with US household wealth against US participation moves the bands by amounts well inside their width (Appendix B).
\item \textbf{Few independent windows.} Five-year returns from a 16-year price history contain roughly three independent observations. That is why the paper leans on levels at measured dates, cross-sectional mortality, and two survey waves rather than on time-series regressions, and why its forecast is phrased so that a single new window can grade it.
\end{enumerate}

\section{Conclusion}

The riddle was that Bitcoin carried three expected returns at once: 32\% from its own tested history, 6--18\% from the CAPM, 14--22\% from its holders' mouths. The resolution is that only one of them was ever an expected return. The CAPM range is the hurdle of an investor who does not set Bitcoin's price. The holders' numbers are the required return of the investor who does, and they move, wave by wave, in the direction and magnitude the marginal-holder condition, Merton's (1987) limited-participation pricing with measured inputs, says they should: with her volatility and her portfolio weight, at ordinary levels of risk aversion. And the 32\% is the smooth curve through a history equal to that required return plus a wedge of unanticipated adoption and winner-selection. The graveyard shows the wedge was never mortality compensation (hazard rises with age; the top-20 jump hazard is 0.3\% a year); the cross-section shows it was never available to the typical asset (the median coin lost money); and Bitcoin's own record shows it turned negative after 2021.

What remains, once the wedge is spent, is an asset priced by observable things: 40-some percent volatility, a holder base diluting toward market weight, a small and datable mortality risk, and a risk-free rate. Those ingredients say 9--22\% a year for 2026--2031. If Bitcoin instead repeats its age-rule 32\%, the market will have found a new surprise we could not measure, and the framework here says exactly where to look for it. Either way, the question that opened the companion paper (compensation for risks that never came, or an asset the market underpriced?) now has an evidence-based answer: on either reading of the early marginal holder, mostly the second, for a decade; neither, anymore.

\appendix

\section{The Curated Elite Death Table}

All 36 deaths among the 219 ever-top-20 assets, with the classification each received. \emph{Type}: jump = sudden, value-destroying collapse priced as a Poisson event; exchange = failure of the issuing exchange; fade = gradual abandonment and loss of listings; contested = flagged by a mechanical rule but the chain still operates. $\delta$ is the approximate fraction of pre-event value recoverable by a holder (last liquid price or successor claim as a share of pre-event value); it is an estimate, not an accounting figure. \emph{Conf.}\ flags classification confidence (H/M/L). Dates are the curated economic death dates; mechanical rules alone would misdate LUNA, UST, and FTT, whose collapses are entered from the public record.

\begin{center}\footnotesize
\begin{longtable}{llllp{5.4cm}l}
\toprule
Asset & Death date & Type & $\delta$ & Note & Conf. \\
\midrule
\endhead
LTL & 2014-12 & fade & $\sim$0 & minor 2013-era coin & M \\
FCT (Factom) & 2021-07 & fade & $\sim$0.05 & sponsor company bankrupt 2020; chain stalled & H \\
LBC (LBRY) & 2021-09 & fade/regulatory & $\sim$0.05 & SEC action 2021; project wound down & H \\
PPT (Populous) & 2021-11 & fade & $\sim$0 & abandoned & M \\
ETP (Metaverse) & 2022-01 & fade & $\sim$0 & abandoned & M \\
NCASH (Nucleus Vision) & 2022-01 & fade & $\sim$0 & delisted majors 2021 & H \\
ZB & 2022-02 & exchange & $\sim$0 & ZB.com halted withdrawals 2022 & H \\
WTC (Waltonchain) & 2022-04 & fade & $\sim$0 & abandoned & H \\
\textbf{UST (TerraUSD)} & \textbf{2022-05-12} & \textbf{jump} & $\sim$0.02 & algorithmic stablecoin depeg & H \\
\textbf{LUNA (Terra)} & \textbf{2022-05-12} & \textbf{jump} & $\sim$0.01 & supply hyperinflation; LUNA2 airdrop gave partial recovery & H \\
QASH (Liquid) & 2022-07 & exchange & $\sim$0 & exchange hacked 2021, absorbed into FTX & H \\
QC (Qcash) & 2022-08 & fade & $\sim$0 & abandoned & M \\
\textbf{FTT (FTX Token)} & \textbf{2022-11-10} & \textbf{jump} & $\sim$0.05 & FTX insolvency & H \\
MITH (Mithril) & 2023-03 & fade & $\sim$0 & serial delistings & H \\
QSP (Quantstamp) & 2023-05 & fade & $\sim$0.02 & token abandoned & M \\
SEELE & 2023-07 & fade & $\sim$0 & abandoned & M \\
VTC (Vertcoin) & 2023-08 & fade & $\sim$0.02 & lost major listings & M \\
MDA (Moeda) & 2023-09 & fade & $\sim$0 & abandoned & M \\
BTS (BitShares) & 2024-06 & fade & $\sim$0.02 & community remnant only & M \\
AE (Aeternity) & 2024-10 & fade & $\sim$0.02 & team wind-down & M \\
BCN (Bytecoin) & 2024-10 & fade & $\sim$0 & long dormant; delisted majors years earlier & H \\
OMG (OMG Network) & 2025-02 & fade & $\sim$0.02 & Binance delisting; development ceased & H \\
KEY (SelfKey) & 2025-04 & fade & $\sim$0 & abandoned & M \\
LOOM (Loom Network) & 2025-05 & fade & $\sim$0 & abandoned & H \\
XEM (NEM) & 2025-05 & fade & $\sim$0.05 & major delistings; earlier opt-in XYM airdrop (2021) & H \\
BTG (Bitcoin Gold) & 2025-06 & fade & $\sim$0.02 & Binance delisting & H \\
LINA (Linear) & 2025-06 & fade & $\sim$0 & project shut down & M \\
PDA (PlayDapp) & 2025-06 & fade & $\sim$0.02 & post-2024-exploit decline (PLA$\to$PDA swap preceded) & M \\
YFII (DFI.Money) & 2025-07 & fade & $\sim$0 & abandoned fork & M \\
FLOW & 2025-09 & \textbf{contested} & n/a & chain operating; $>$99\% below ATH & H \\
MXC & 2025-09 & fade & $\sim$0.02 & rebranding/decline & L \\
STMX (StormX) & 2025-10 & fade & $\sim$0 & abandoned & M \\
NANO & 2025-12 & \textbf{contested} & n/a & active community; lost major listings & M \\
ICP (Internet Computer) & 2026-01 & \textbf{contested} & n/a & alive, top-100; ``death'' is an artifact of the listing-day price-spike ATH & H \\
HIFI & 2026-01 & fade & $\sim$0 & delisted & L \\
SUSHI (SushiSwap) & 2026-03 & \textbf{contested} & n/a & protocol operating; $>$99\% below ATH & H \\
\bottomrule
\end{longtable}
\end{center}

\textbf{Headline sensitivities.} All 36 deaths: hazard 3.54\%/yr (CI 2.48--4.90). Excluding the four contested cases: 3.15\%/yr (CI 2.15--4.44). Jump-collapses only (LUNA, UST, FTT): 0.30\%/yr (CI 0.06--0.86). In the ever-top-50 universe the jump set adds SRM (Serum) and ANC (Anchor), both 2022 casualties of the same two cascades, for 0.23\%/yr.

\textbf{Migrations excluded as deaths} (holders received successor value): MATIC$\to$POL, EOS$\to$Vaulta (A), FTM$\to$S (Sonic), GNT$\to$GLM, RNDR$\to$RENDER, BTT redenomination, NU$\to$T (Threshold merger), STRAT$\to$STRAX, MAID$\to$eMAID, HT$\to$HTX, REP$\to$REPv2, AGIX$\to$FET (ASI merger), XZC$\to$FIRO, TOMO$\to$VIC, NPXS$\to$PUNDIX, BNX redenomination, and DGD, whose dissolution refunded the ETH treasury to holders: a death with $\delta$ near 1, the mirror image of LUNA.

\section{Regressions and Robustness}

\textbf{B.1 The Lindy test.} Poisson regression of monthly death indicators on log age, offset $\log(1/12)$, standard errors clustered by asset. Lindy (hazard $= k/\text{age}$) requires a log-age coefficient of $-1$.

\begin{table}[ht]\centering\small
\begin{tabular}{lccc}
\toprule
Death definition & Deaths & $\gamma_{\text{age}}$, no FE & $\gamma_{\text{age}}$, year FE \\
\midrule
Strict (feed termination) & 234 & $+0.56$ (0.07) & $+0.44$ (0.06) \\
Broad (feed + price) & 263 & $+0.69$ (0.07) & $+0.56$ (0.07) \\
Abandonment (all mechanical) & 588 & -- & $+0.51$ (0.04) \\
\textbf{Curated (final)} & \textbf{520} & -- & $\boldsymbol{+0.42}$ \textbf{(0.04)} \\
Curated, 2025 feed-wall removed & -- & -- & $+0.65$ (0.08) \\
Curated, birth-cohort FE in place of year FE & 520 & -- & $+1.62$ (0.10) \\
Elite (ever top-20) subsample & 36 & -- & $+0.68$ (0.21) \\
\bottomrule
\end{tabular}
\end{table}

The log-survival--log-age slope over ages 1--10 runs from $-0.21$ (strict) to $-0.49$ (curated), i.e.\ a cumulative Pareto constant $\hat{k} \approx 0.2$--$0.5$ across definitions, an order of magnitude below the 5.6 the age rule would require. Two further robustness checks address the standard identification worry that an age effect might be a disguised cohort effect (age, cohort, and period are collinear, so only two can enter at once): replacing calendar-period fixed effects with birth-cohort fixed effects \emph{raises} the age elasticity to $+1.62$ (s.e.\ 0.10), so the rejection of Lindy is not an artifact of the period controls, and the period-FE $+0.42$ is the conservative figure. Parametrically, a Weibull fit to the curated durations gives shape $\rho = 1.57$ (95\% CI 1.47--1.67), where $\rho > 1$ means a hazard rising with age and Lindy requires $\rho < 1$; a log-logistic fit (slightly preferred by AIC, 3,342 vs 3,371) implies a hump-shaped hazard that rises through the first decade of life. The implied median lifetime of a core-universe cryptoasset is 6.4 years under both the Weibull and the Kaplan--Meier estimates.

\textbf{B.2 The concentration quantity is not a timing signal.} Forward $h$-year log drift regressed on the real-time concentration quantity $x = \sigma^2 \cdot (\text{cap}/\text{participating wealth})$ and the age-rule regressor $z = \ln(1+5/A)/5$; heteroskedasticity- and autocorrelation-robust (HAC) standard errors with $h+6$ lags; overlapping monthly windows from 2012 (independent observations $\approx$ 2--4, so read $t$-statistics as descriptive).

\begin{table}[ht]\centering\small
\begin{tabular}{lccccc}
\toprule
Horizon & $x$ alone & $z$ alone & Horse race: $x$ & Horse race: $z$ & $R^2$ (joint) \\
\midrule
3y & $-45.1$ ($t$ $-5.9$) & $+6.0$ ($t$ 5.0) & $-39.1$ & $+5.1$ & 0.67 \\
5y & $-33.9$ ($t$ $-4.1$) & $+7.0$ ($t$ 13.5) & $-30.1$ & $+6.5$ & 0.81 \\
\bottomrule
\end{tabular}
\end{table}

The concentration quantity enters with the \emph{wrong sign} as a predictor (market capitalization in its numerator makes it spike at mania tops) while the age regressor replicates the companion paper's coefficient ($\approx$ 5--7, the power-law exponent). This is why the paper treats the marginal-holder condition as a date-by-date equilibrium restriction and never as a forecasting variable; conversely, it is also why the age rule's in-sample fit should not be read as evidence of a priced mechanism.

\textbf{B.3 The companion paper's own metric, applied to a naive version of this model.} A reader might ask the required-return path to pass the backtest that selected the age rule. It does not, and should not: evaluated as a 5-year forecast over 2012--2021, the survey-anchored required-return path is beaten in 82--88\% of windows by 111--125 points a year, numerically similar to the 7--21\% institutional rates the companion paper rejects. That is the wedge, measured a third way: over the sample in which the age rule was fitted, realized returns exceeded any defensible contemporaneous required return by roughly 100 points a year on average. A rate can be the correct required return and a poor forecast of a sample dominated by surprises; the backtest cannot distinguish an unbiased forecast from a reduced form that happens to track surprise decay.

\textbf{B.4 Pairing robustness.} Re-running the market-clearing back-cast with US household net worth (Federal Reserve Z.1, $\approx$ \$80T in 2013 to $\approx$ \$170T in 2026) against the US participation series, instead of global wealth: the 2021 required band becomes 22--39\% (vs 22--58\%), the 2026 band 13--22\% (vs 9--22\%, which also reflects wider $w$ bounds). All qualitative statements survive; the 2026 forecast band is quoted at its widest.

\textbf{B.5 Death-definition sensitivities.} Volume-collapse deaths require an absolute floor (\$250k/day trailing volume) in addition to the 1\%-of-peak rule; without the floor, 2021-mania base effects falsely kill operating mid-caps. The elite hazard is 3.54\% with all curated deaths, 3.15\% excluding the four contested zombies, and 2.36\% (CI 1.51--3.51\%) treating all twelve 2025 deaths as censoring; the Lindy rejection holds under every variant (B.1).

\section*{References}
\begin{list}{}{\setlength{\itemindent}{-1.5em}\setlength{\leftmargin}{1.5em}\setlength{\itemsep}{2pt}\setlength{\parsep}{0pt}}
\item Bekaert, G., and C. R. Harvey. 2000. ``Foreign Speculators and Emerging Equity Markets.'' \emph{Journal of Finance} 55(2): 565--613.
\item Biais, B., C. Bisi\`ere, M. Bouvard, C. Casamatta, and A. J. Menkveld. 2023. ``Equilibrium Bitcoin Pricing.'' \emph{Journal of Finance} 78(2): 967--1014.
\item Black, F., and R. Litterman. 1992. ``Global Portfolio Optimization.'' \emph{Financial Analysts Journal} 48(5): 28--43.
\item Brown, S., W. Goetzmann, and S. Ross. 1995. ``Survival.'' \emph{Journal of Finance} 50(3): 853--873.
\item Burger, H. C. 2019. ``Bitcoin's Natural Long-Term Power-Law Corridor of Growth.'' Coinmonks (Medium).
\item Candia, B., O. Coibion, Y. Gorodnichenko, and M. Weber. 2026. ``Do You Even Crypto, Bro? Cryptocurrencies in Household Finance.'' Federal Reserve Bank of Cleveland Working Paper 26-16.
\item Coin Metrics. 2026. ``Community Network Data.'' GitHub repository coinmetrics/data. Accessed November 2026.
\item Cong, L. W., Y. Li, and N. Wang. 2021. ``Tokenomics: Dynamic Adoption and Valuation.'' \emph{Review of Financial Studies} 34(3): 1105--1155.
\item Deloitte. ``Quantum Computers and the Bitcoin Blockchain.'' Deloitte Netherlands analysis (standing estimate of quantum-vulnerable supply).
\item Deribit. DVOL implied-volatility index readings as reported January 2026 (level 40; one-year average 55).
\item Feder, A., N. Gandal, J. T. Hamrick, and T. Moore. 2018. ``The Rise and Fall of Cryptocurrency Coins and Tokens.'' \emph{Decisions in Economics and Finance} 41: 195--222.
\item Federal Reserve Board. 2022--2024. \emph{Survey of Household Economics and Decisionmaking (SHED)}, crypto-holding items; Federal Reserve Banks of Boston/Atlanta/San Francisco, \emph{Survey and Diary of Consumer Payment Choice}, bitcoin-ownership items.
\item Giglio, S., M. Maggiori, J. Stroebel, and S. Utkus. 2021. ``Five Facts about Beliefs and Portfolios.'' \emph{American Economic Review} 111(5): 1481--1522.
\item Grobys, K., and N. Sapkota. 2020. ``Predicting Cryptocurrency Defaults.'' \emph{Applied Economics} 52(46): 5060--5076.
\item Henry, P. B. 2000. ``Stock Market Liberalization, Economic Growth, and the Cost of Capital.'' \emph{Journal of Finance} 55(2): 529--564.
\item Liu, Y., and A. Tsyvinski. 2021. ``Risks and Returns of Cryptocurrency.'' \emph{Review of Financial Studies} 34(6): 2689--2727.
\item Makarov, I., and A. Schoar. 2020. ``Trading and Arbitrage in Cryptocurrency Markets.'' \emph{Journal of Financial Economics} 135(2): 293--319.
\item Martin, I. 2017. ``What Is the Expected Return on the Market?'' \emph{Quarterly Journal of Economics} 132(1): 367--433.
\item Merton, R. C. 1969. ``Lifetime Portfolio Selection under Uncertainty: The Continuous-Time Case.'' \emph{Review of Economics and Statistics} 51(3): 247--257.
\item Merton, R. C. 1976. ``Option Pricing When Underlying Stock Returns Are Discontinuous.'' \emph{Journal of Financial Economics} 3(1--2): 125--144.
\item Merton, R. C. 1987. ``A Simple Model of Capital Market Equilibrium with Incomplete Information.'' \emph{Journal of Finance} 42(3): 483--510.
\item OeNB (Oesterreichische Nationalbank). 2025. ``Crypto Assets in Austria: Evidence from HFCS Wave 5.'' \emph{OeNB Bulletin} Q4/25.
\item Pagnotta, E. S. 2022. ``Decentralizing Money: Bitcoin Prices and Blockchain Security.'' \emph{Review of Financial Studies} 35(2): 866--907.
\item Papadogiannis, A. 2025. ``Economic Scarcity as Value: A Universal Framework for Blockchain and Beyond.'' Working paper, under revision at \emph{IET Blockchain}.
\item Papadogiannis, A. 2026a. ``A Dual-Risk Theory of Valuation Separating Market-Based Discounting from Expectational Distortion.'' \emph{Discover Analytics} 4(1): 19.
\item Papadogiannis, A. 2026b. ``From Unified Discounting to Dual-Risk: A New Lens on DCF Valuation.'' \emph{Finance Research Open} 2: 100096.
\item Papadogiannis, A. 2026c. ``Bitcoin's Expected Return: An Age Rule That Beats Current Practice.'' Working paper, submitted to the \emph{Journal of Alternative Investments}.
\item Pew Research Center. 2021--2024. Shares of US adults who have invested in, traded, or used cryptocurrency.
\item ``Risk Premia in the Bitcoin Market.'' 2024. arXiv:2410.15195 (Deribit options, July 2017--December 2022).
\item UBS. Annual. \emph{Global Wealth Report} (global household wealth aggregates).
\end{list}

\vspace{1em}
\noindent\textbf{Data and code availability:} the full replication package (panel-construction and survival scripts, the curated asset-level panel, all result tables, and every exhibit) is maintained at \url{https://github.com/apapadogiannis/bitcoin-required-return}. The pipeline ends with an audit script that recomputes and checks every quantitative claim in this paper. Underlying market data are the public Coin Metrics community dataset (snapshot ending May 24, 2026), which the package downloads from its source rather than redistributing.

\vspace{1em}
\noindent\textit{Disclaimer: This research was conducted independently and without funding, review, or endorsement. The views expressed are solely those of the author and do not represent any institutional position. This paper is for informational and academic purposes only and does not constitute financial, investment, tax, or legal advice, nor a recommendation to buy or sell any financial instrument.}

\end{document}
